\documentclass[11pt]{article}
\usepackage[margin=1in]{geometry}
\usepackage{amsmath,amssymb}
\usepackage[hidelinks]{hyperref}

\newcommand{\APEP}{\operatorname{APEP}}
\newcommand{\ext}{\operatorname{ext}}
\newcommand{\wstar}{w^*}

\begin{document}

\section*{Human Review: Interior APEPs in \texorpdfstring{$\ell_p$}{ell-p}-sums}

\noindent Original automatic proof: \texttt{solution\_packet.pdf}

\noindent Checked date: 18 June 2026

\noindent Status: Manually checked.

\noindent Verdict: The proof appears correct.

\section*{Human Comments}

The proposed characterization appears correct and quite natural.  The source
paper gives the general criterion
\[
  x\in \operatorname{ape}(C)
  \quad\Longleftrightarrow\quad
  x\in \overline{\ext(C^{\wstar})}^{\,\wstar},
\]
and Proposition 4.9 characterizes the norm-one APEPs in \(\ell_p\)-sums.
Question 8.2 asks for the remaining interior case.  This packet supplies a
full characterization by introducing the admissible coordinate radii.

\section*{Main Idea}

For each coordinate, set
\[
  C_i=\overline{\ext B_{X_i^{**}}}^{\,\wstar}.
\]
For \(u\in X_i\), a radius \(r\in[0,1]\) is admissible if
\[
  u\in r C_i.
\]
The least admissible radius \(\rho_i(u)\) measures how much \(p\)-mass is
needed in the \(i\)-th coordinate to realize \(u\) as a weak-star limit of
scaled coordinate extreme points.

With this notation, the characterization is very natural.  In the infinite
summand case, \(x=(x_i)\) is an interior APEP precisely when
\[
  \sum_i \rho_i(x_i)^p\leq 1.
\]
The remaining \(p\)-mass can be placed in a coordinate outside any finite
weak-star test.  In the finite summand case, there is no outside coordinate,
so one needs admissible radii \(r_i\in R_i(x_i)\) with exact equality
\[
  \sum_i r_i^p=1.
\]

\section*{Proof Check}

The necessity argument is sound.  If \(x\) is an APEP, Theorem 3.8 gives a
net of extreme points
\[
  y^\alpha=(y_i^\alpha)\in \ext B_{X^{**}}
\]
weak-star converging to \(x\).  Coordinatewise weak-star convergence follows
from the coordinate embeddings \(X_i^*\hookrightarrow X^*\).  For each finite
set of coordinates, one passes to a subnet so that
\[
  \|y_i^\alpha\|\to r_i.
\]
If \(r_i=0\), then \(x_i=0\), hence \(0\in R_i(x_i)\).  If \(r_i>0\), then
\[
  \frac{y_i^\alpha}{\|y_i^\alpha\|}
  \in \ext B_{X_i^{**}},
\]
and after passing to a further subnet one obtains \(x_i=r_i e_i\) with
\(e_i\in C_i\).  Thus \(r_i\in R_i(x_i)\) and
\(\rho_i(x_i)\leq r_i\).  Since every \(y^\alpha\) has norm one, finite
partial sums satisfy
\[
  \sum_{i\in F}\rho_i(x_i)^p\leq 1.
\]
Taking the supremum over finite \(F\) gives the required inequality in the
infinite case, while taking \(F\) to be the finite support gives the equality
condition in the finite case.

The sufficiency argument is also sound.  In the finite case, admissible radii
with \(p\)-sum equal to one allow one to approximate each coordinate by genuine
coordinate extreme points and thereby build global extreme points of
\(B_{X^{**}}\) weak-star close to \(x\).  In the infinite case, finite-support
approximation in
\[
  X^*=\left(\bigoplus_i X_i^*\right)_q
\]
reduces a weak-star neighbourhood to finitely many coordinates.  One realizes
the required coordinates using the least admissible radii, and if there is
leftover mass, places it in an unused coordinate.  The extreme-point formula
for \(\ell_p\)-sums then gives an extreme point of \(B_{X^{**}}\) inside the
chosen neighbourhood.

\section*{Review Outcome}

Approve this packet as an apparently correct full solution to Question 8.2.
The characterization is not simply copied from the paper: the paper supplies
the general weak-star closure criterion and the norm-one \(\ell_p\)-sum case,
while this argument extends the analysis to interior points through the
admissible-radius formulation.

\end{document}
